% !TEX root = main.tex

% ======================================================================
% Template Showcase
% Demonstrates all available environments and macros in preamble.tex
% ======================================================================

% --- Slide 1: Overview ---
\blueheader
\begin{frame}{Template Showcase: Overview}
This presentation demonstrates the building blocks available in this template.
\medskip
\begin{itemize}
  \item \textbf{Header colors}: \texttt{\textbackslash blueheader}, \texttt{\textbackslash redheader}, \texttt{\textbackslash greenheader}, \texttt{\textbackslash cyanheader}, \texttt{\textbackslash grayheader}
  \item \textbf{Theorem boxes}: \texttt{bluebox}, \texttt{redbox}, \texttt{greenbox}, \texttt{cyanbox}, \texttt{graybox}
  \item \textbf{Focus boxes}: \texttt{bluefocus}, \texttt{redfocus}, \texttt{greenfocus}, \texttt{cyanfocus}, \texttt{grayfocus}
  \item \textbf{Margin macros}: \texttt{\textbackslash warn}, \texttt{\textbackslash info}, \texttt{\textbackslash warnCustomMsg}, \texttt{\textbackslash infoCustomMsg}
  \item \textbf{Math macros}: \texttt{\textbackslash norm\{\}}, \texttt{\textbackslash abs\{\}}
\end{itemize}
\end{frame}

% --- Slide 2: Blue ---
\blueheader
\begin{frame}{Template Showcase: Blue}
\begin{columns}[T]
  \begin{column}{0.48\textwidth}
    \begin{bluebox*}{Definition}
      A \textbf{set} $S$ is called \emph{finite} if there exists a bijection $S \to \{1,\ldots,n\}$ for some $n \in \mathbb{N}$.
    \end{bluebox*}
    \medskip
    \begin{bluebox}{Definition (with label)}{def:finite}
      Same box, but numbered and referenceable via \texttt{def:finite}.
    \end{bluebox}
  \end{column}
  \begin{column}{0.48\textwidth}
    \begin{bluefocus}
      \textbf{bluefocus} --- a titled-free highlight box. Use it to draw attention to a key formula or remark without a formal theorem label.
      \[
        e^{i\pi} + 1 = 0
      \]
    \end{bluefocus}
  \end{column}
\end{columns}
\end{frame}

% --- Slide 3: Red ---
\redheader
\begin{frame}{Template Showcase: Red}
\begin{columns}[T]
  \begin{column}{0.48\textwidth}
    \begin{redbox*}{Theorem}
      For all $a, b \in \mathbb{R}$,
      \[
        (a + b)^2 = a^2 + 2ab + b^2.
      \]
    \end{redbox*}
    \medskip
    \begin{redbox}{Theorem (numbered)}{thm:binomial}
      Same box with a counter shared across all box types.
    \end{redbox}
  \end{column}
  \begin{column}{0.48\textwidth}
    \begin{redfocus}
      \textbf{redfocus} --- use for warnings, important notes, or anything that demands immediate attention.

      \medskip
      Avoid using \emph{red} purely decoratively; reserve it for critical content.
    \end{redfocus}
  \end{column}
\end{columns}
\end{frame}

% --- Slide 4: Green ---
\greenheader
\begin{frame}{Template Showcase: Green}
\begin{columns}[T]
  \begin{column}{0.48\textwidth}
    \begin{greenbox*}{Example}
      Let $f(x) = x^2$. Then $f'(x) = 2x$, so the slope at $x = 3$ is
      \[
        f'(3) = 6.
      \]
    \end{greenbox*}
    \medskip
    \begin{greenbox}{Example (numbered)}{ex:deriv}
      Same box, numbered.
    \end{greenbox}
  \end{column}
  \begin{column}{0.48\textwidth}
    \begin{greenfocus}
      \textbf{greenfocus} --- ideal for worked examples, intuition, or positive reinforcement of a concept just introduced.
    \end{greenfocus}
  \end{column}
\end{columns}
\end{frame}

% --- Slide 5: Cyan ---
\cyanheader
\begin{frame}{Template Showcase: Cyan}
\begin{columns}[T]
  \begin{column}{0.48\textwidth}
    \begin{cyanbox*}{Exercise}
      Compute $\displaystyle\int_0^1 x^2\,dx$ and verify that the answer equals $\tfrac{1}{3}$.
    \end{cyanbox*}
    \medskip
    \begin{cyanbox}{Exercise (numbered)}{ex:integral}
      Same box, numbered.
    \end{cyanbox}
  \end{column}
  \begin{column}{0.48\textwidth}
    \begin{cyanfocus}
      \textbf{cyanfocus} --- well-suited for exercises, tasks, or student prompts that do not need a formal label.
    \end{cyanfocus}
  \end{column}
\end{columns}
\end{frame}

% --- Slide 6: Gray ---
\grayheader
\begin{frame}{Template Showcase: Gray}
\begin{columns}[T]
  \begin{column}{0.48\textwidth}
    \begin{graybox*}{Exploration}
      Pick any two vectors $u, v \in \mathbb{R}^2$. Compute $\norm{u+v}$ and compare it with $\norm{u} + \norm{v}$. What do you notice?
    \end{graybox*}
    \medskip
    \begin{graybox}{Exploration (numbered)}{ex:triangle}
      Same box, numbered.
    \end{graybox}
  \end{column}
  \begin{column}{0.48\textwidth}
    \begin{grayfocus}
      \textbf{grayfocus} --- suited for open-ended exploration, background remarks, or supplementary context.
    \end{grayfocus}
  \end{column}
\end{columns}
\end{frame}

% --- Slide 7: Macros ---
\blueheader
\begin{frame}{Template Showcase: Macros}

\textbf{Warning and info macros:}
\medskip
\begin{itemize}
  \item \warn
  \item \info
  \item \warnCustomMsg[myblue]{Custom warning in any color}
  \item \infoCustomMsg[mygreen]{Custom info in any color}
\end{itemize}

\medskip
\textbf{Math macros:}
\medskip
\begin{itemize}
  \item Norm: $\norm{v} = \sqrt{\langle v, v \rangle}$
  \item Absolute value: $\abs{x - y} \leq \abs{x} + \abs{y}$
  \item Combined: $\abs{\langle u, v \rangle} \leq \norm{u}\,\norm{v}$ \quad (Cauchy--Schwarz)
\end{itemize}

\end{frame}
